\section{Purpose and scope}

\emph{Fractal Fusion Scales} is A$\Omega$D Manual II: the single application manual for the fusion ladder up the octaves. It compacts application material into one scale ladder because the operative object is one repeated AFC/A$\Omega$D procedure rather than multiple domain-specific calculi.

Manual II demonstrates one fusion-ladder machinery across declared octaves:
\[
(\zeta,R_\zeta,\partial R_\zeta,\omega)
\longrightarrow
\text{raw AFC/D.E.C. trace}
\longrightarrow
\text{trace detection}
\longrightarrow
\text{frozen rows or target packets}
\longrightarrow
\Obs.
\]
The octave changes the observable dictionary, fixture data, target packet, and comparison source. The calculus procedure does not change.

\subsection{What is new in Manual II}

Manual II is the first-class source subtree \texttt{manual-2/}. It carries deterministic fixture and target-packet lanes:
\begin{itemize}[leftmargin=2em]
  \item an elementary octave using the declared $3{:}3{:}6$ and $3{:}4{:}6$ fusion-scale ladders against a frozen 118-element registry fixture;
  \item a molecular octave using exact formula accounting for small precursor and closure fixtures;
  \item a molecular-data target scaffold for PubChem locator packets, RDKit graph/canonicalization fixture rows, component vectors, route units, controls, and link-rule schemas;
  \item a bio-chain octave using nucleotide, dinucleotide, RNA/DNA chain, peptide, and protein-chain candidate closures;
  \item a protein target lane for UniProt sequence locators, PDB/mmCIF experimental-structure locators, AlphaFold predicted-structure comparator locators, target-normalization rows, value-map quarantine rows, and target-leakage guards.
\end{itemize}

\subsection{Claim discipline}

Status. Manual II records a reproducible fusion-ladder protocol and exact fixture, target-packet, freeze-row, and residual-audit lanes. External references enter only after the relevant rows are frozen or as target locators for future comparison gates.

\aodnoteblock{Guardrail}{The elementary rows are registry-coordinate and scale-comparison rows. The molecular and bio-chain rows are formula-closure rows. PubChem, RDKit, UniProt, PDB/mmCIF, and AlphaFold rows are target or comparison lanes. None of those external sources becomes a premise of the raw AFC/A$\Omega$D D.E.C. run.}

\subsection{Octave convention}

An \emph{octave} is a declared level/window/boundary/readout package:
\[
\mathcal O_\ell=(\zeta_\ell,R_\ell,\partial R_\ell,\omega_\ell,\Detect_\ell,\Pi_{\mathrm{obs},\ell},\Data_\ell).
\]
The ladder is fractal in the limited operational sense that the same finite boundary procedure is rerun under different readout dictionaries. Manual II now uses these lane names:
\begin{center}
\begin{tabularx}{\linewidth}{@{}lX@{}}
\toprule
Lane & Readout or target dictionary \\
\midrule
Elementary & atomic-number registry coordinate $Z$, ladder depth, retained capacity, residual-$Z$, scale-comparison label \\
Molecular & component formula, fusion closure, shedding vector, CHNOPS residual ledger, PubChem/RDKit target packet \\
Bio-chain & sequence, nucleotide or residue closure, chain formula, water-shedding count, RDKit graph lane \\
Protein target & UniProt locator, PDB/mmCIF locator, AlphaFold comparator locator, contact-map hash, distance-matrix hash, leakage guard, deferred folding/contact-map value map \\
\bottomrule
\end{tabularx}
\end{center}
